\pagestyle{fancy}
\fancyfoot{}
\fancypagestyle{firstpage}{
\fancyfoot[RO, RE]{\fontsize{15}{16}\selectfont 36}{\rightmark}
\fancyhead[C]{\fontsize{12}\selectfont \textbf{КВАНТ} · \textsf{2 0 1 7 / № 3}}
}

\definecolor{mypurple}{RGB}{165,164,208} % определия свой цвет
\fancypagestyle{secondpage}{
\begin{tikzpicture}[remember picture,overlay]
\fill[mypurple] (-0.5,0.7) rectangle (50,20)
\end{tikzpicture}
\fancyhead[C]{\Large\textbf{Физический факультатив}}
\fancyhead[R]{\Large37}}
\renewcommand{\headrulewidth}{0pt}

\setlength{\columnsep}{1 cm}

\thispagestyle{firstpage}
\fontsize{12}{14}\selectfont
\begin{multicols}{2}
\hspace{6cm}Таблица 1.


\resizebox{\columnwidth}{!}{

\begin{tabular}{|c|c|c|c|c|c|c|c|c|}
    \hline
    \rowcolor{gray!50}Xe & Ar & Kr & Ne & Rn & CBr$_4$ & C$_6$H$_{12}$ & H$_2$S & N$_2$ \\
    \hline
    1,80 & 1,80 & 1,80 & 1,82 & 1,87 & 1,95 & 1,98 & 1,99 & 2,00 \\
    \hline
    \rowcolor{gray!50}C$_6$H$_6$ & Br$_2$ & CH$_4$ & I$_2$ & CCl$_4$ & CrH$_{14}$ & H$_2$O & H$_2$ & Cl$_2$ \\
    \hline
    2,02 & 2,08 & 2,11 & 2,14 & 2,22 & 2,32 & 2,37 & 2,39 & 2,42 \\
    \hline
    \rowcolor{gray!50}CF$_4$ & C$_8$H$_{18}$ & F$_2$ & C$_9$H$_6$ & O$_2$ & CsH$_{10}$ & O=C=Me$_2$ & P$_2$ & Pt \\
    \hline
    2,55 & 2,63 & 2,70 & 2,73 & 2,86 & 2,86 & 2,86 & 3,05 & 3,15 \\
    \hline
    \rowcolor{gray!50}He & W & S$_2$ & C$_2$H$_6$ & Ag & Au & Zn & Mo & Cu \\
    \hline
    3,25 & 3,25 & 3,35 & 3,41 & 3,49 & 3,61 & 3,74 & 3,85 & 3,97 \\
    \hline
    \rowcolor{gray!50}Cd & Zr & Pb & Rb & Na & K & Cs & Hg & Li \\
    \hline
    4,18 & 4,21 & 6,61 & 6,74 & 6,75 & 6,78 & \multicolumn{3}{r|}{P3109} \\
    \hline
\end{tabular}

}

цепочек атомов (молекул), которые могут
изгибаться, и между цепочками имеются
промежутки. Если для щелочных металлов
$Z_2 \approx 7$, то это означает, что значительная
часть вещества в среднем по времени пребывает в составе неких трубок-цилиндров-струй
с заполнением. На каждую молекулу, находящуюся на поверхности такой трубки, приходится по шесть таких же соседок, живущих на поверхности, и одна – внутри трубки. А в промежутках между трубками –
пустота. Это можно представить как структуру пены с тонкими стенками и толстыми
участками, на которых стенки соединяются
друг с другом под неким углом в пространстве, или как модель кристаллической решетки, у которой стерженьки, соединяющие
узлы решетки, являются теми самыми цилиндрами-трубками-струями.
\indent Каждая молекула, входящая в состав конденсированного вещества, при невысоких
внешних давлениях и при температуре значительно меньше критической занимает
объем, примерно равный $D^3$ D . А кинетическая энергия поступательного движения молекулы равна $E_{кин}=3kT/2$. Таким образом, давление, связанное с тепловым движением, которое молекула оказывает на стенки своей «ячейки/клетки», равно $(2/3)E_{кин}/D^3=kT/D^3=RT\rho/M$. Когда эта
величина сравняется с собственным давлением вещества при числе соседок $Z_2$, вещество не сможет находиться в конденсированном состоянии, т.е. это равенство дает
еще один критерий нахождения критической температуры: $RT\rho/M=2Z_2L\rho/(ZM)$. Отсюда следует $$T_{кр2}=\frac{2Z_2L}{Z_0R}$$ Если разделить критическую температуру $T_{кр2}$ на критическую температуру $T_{кр1}$ , то получится безразмерное число: $T_{кр2}/T_{кр1}\approx i-2$. Заметим, что при i = 3, т.е. для веществ с одноатомными молекулами, это отношение равно единице, значит, разные подходы дают один и тот же результат для критической температуры.
\section*{\centering\large\textit{\textcolor{red}{Поверхностная энергия и
коэффициент поверхностного натяжения}}}


Молекулы конденсированного вещества,
которые живут на границе раздела с паром
этого же вещества, имеют меньшее количество ближайших соседок $Z_3$, чем молекулы,
живущие внутри объема и имеющие число
соседок $Z$. Поэтому потенциальная яма, в
которой находится каждая молекула на поверхности, имеет меньшую глубину.
\indent При невысоких температурах, когда взаимодействием молекулы на поверхности с молекулами пара можно пренебречь, у каждой молекулы на поверхности в среднем число соседок меньше на определенную долю от максимального числа соседок: $\Delta Z=Z-Z_3$3. Например, если $Z = 12$ при плотной упаковке
шариков, то у молекул на плоской поверхности соседок всего $Z_3=9$. Следовательно потенциальная энергия в положении равновесия у таких молекул равна $-9U_0$, а избыточная энергия равна $3U_0=ZU_0/4$. Каждая молекула на поверхности занимает площадь, которая по порядку величины равна $D^2$. Для самой плотной упаковки шариков эта
площадь составляет $\sqrt{3}D^2/2$. Поэтому избыточная энергия, приходящаяся на единицу площади, равна примерно $$\sigma_0=\left(\frac{ZU_0}{2\sqrt{3}D^2}\right)$$ Это и есть коэффициент поверхностного натяжения при невысоких температурах. По мере роста температуры давление и плотность насыщенного пара растут, и молекулы пара создают для каждой молекулы,



живущей на поверхности конденсированного вещества, дополнительную глубину потенциальной ямы. Кроме того, с ростом температуры при внешнем давлении, равном давлению насыщенного пара вещества, убывает плотность вещества в конденсированном 
% \newpage
состоянии. Из этого следует, что коэффициент поверхностного натяжения $\sigma$ убывает с температурой и при критической температуре обращается в ноль, так как при
этой температуре пропадает различие плотностей конденсированного вещества и его насыщенного пара. Коэффициент поверхностного натяжения составляет определенную долю от теплоты испарения при данной температуре, а теплота испарения с ростом температуры убывает. Качественно описать поведение коэффициента поверхностного натяжения можно, например, такой формулой: $$\sigma(T)=\sigma_0 \left( 1-\frac{T}{T_{кр}}\right)$$ Данные из справочника физических величин дают основание для выбора именно
такой зависимости (рис. 1). Если нужно

\includegraphics{photo1.png}

\caption{\textit{Рис. 1. Поверхностное натяжение сжиженных
газов}}
\end{multicols}
\thispagestyle{secondpage}

% \section{Introduction}

